\section{Overlays, logos y texto dinámico}

Voctomix 2.0 incorpora un sistema completo de superposición de elementos gráficos
sobre la señal de vídeo: logos corporativos, overlays PNG con canal alfa y texto
dinámico renderizado en tiempo real. Todos estos elementos se gestionan desde la
interfaz gráfica sin interrumpir el pipeline de procesamiento.

\subsection{Logos}

El sistema soporta hasta dos logos simultáneos, superpuestos en posiciones fijas del
encuadre definidas mediante la configuración INI:

\begin{itemize}
  \item \textbf{logo1}: posicionado en la esquina inferior derecha del fotograma.
  \item \textbf{logo2}: posicionado en la esquina superior izquierda del fotograma.
\end{itemize}

Los ficheros PNG de los logos se definen en las secciones \texttt{[logo1]} y
\texttt{[logo2]} del fichero de configuración, permitiendo adaptar la identidad visual
del evento sin modificar el código. La visibilidad de cada logo se puede activar o
desactivar individualmente desde la interfaz gráfica.

\subsection{Overlays PNG y texto dinámico}

Los overlays son imágenes PNG con canal alfa que se superponen sobre la señal de vídeo
mediante el elemento GStreamer \texttt{gdkpixbufoverlay}. El módulo
\texttt{voctocore/lib/overlay.py} gestiona el ciclo de vida completo de un overlay:
carga del fichero PNG, fundido de entrada (\textit{blend-in}), mantenimiento en
pantalla y fundido de salida (\textit{blend-out}). El tiempo de fundido se controla
mediante el parámetro \texttt{blend-time} en la sección \texttt{[overlay]} del
fichero INI (en milisegundos).

Sobre la plantilla PNG activa se puede renderizar texto dinámico (nombre del ponente,
título de la sesión) utilizando la biblioteca Cairo. El texto se introduce en la
interfaz gráfica y se renderiza en tiempo real generando la imagen PNG resultante,
que se envía al núcleo mediante el comando \texttt{set\_overlay}. Esto permite a los
operadores actualizar los rótulos de ponente durante la emisión sin preparación previa
de imágenes.

\subsection{Auto-off de overlays}

La funcionalidad \textit{auto-off} resuelve un problema frecuente en producción en
directo: el operador realiza un corte de cámara olvidando retirar el overlay activo,
que queda superpuesto sobre la nueva fuente de forma inapropiada (por ejemplo, el
nombre de un ponente apareciendo sobre la imagen de otro).

Cuando la opción \texttt{auto-off = true} está activa en la configuración, el módulo
\texttt{overlay.py} registra un \textit{listener} sobre el evento \texttt{video\_status}.
Cada vez que la fuente de vídeo activa cambia, el listener inicia automáticamente la
secuencia de fundido de salida del overlay activo. Si hay texto dinámico superpuesto,
éste se elimina junto al overlay gráfico, garantizando que ningún elemento de
rotulación quede huérfano en pantalla.

El comportamiento puede configurarse mediante:
\begin{itemize}
  \item \texttt{auto-off = true/false} --- activa o desactiva la funcionalidad.
  \item \texttt{user-auto-off = true/false} --- expone un botón de alternancia en la
        GUI para que el operador pueda activar/desactivar el auto-off en tiempo real.
  \item \texttt{blend-time = <ms>} --- tiempo de fundido de salida en milisegundos.
\end{itemize}

% --- Figura sugerida ---
% \begin{figure}[H]
%   \centering
%   \includegraphics[width=0.85\textwidth]{figuras/overlay_ejemplo.png}
%   \caption{Fotograma con overlay activo (rótulo de ponente) sobre imagen de cámara.}
%   \label{fig:overlay_ejemplo}
% \end{figure}
